\section{Three Flavors of Mutex}

\newcommand{\spectriple}[3]{%
  \begingroup
  \setlength{\fboxsep}{1.2pt}%
  \setlength{\fboxrule}{0.3pt}%
  \boxed{\begin{gathered}
    \left\{\,#1\,\right\} \\[-1pt]
    #2 \\[-1pt]
    \left\{\,#3\,\right\}
  \end{gathered}}%
  \endgroup}

% Keep neighboring triples close while the boxes provide their visual division.
\renewcommand{\MathparLineskip}{\lineskiplimit=.55em\lineskip=.55em plus .1em}


\subsection{\texorpdfstring{\lstinline{std::mutex}}{std::mutex}}

\begin{figure}[htb]
\begin{lstlisting}
std::atomic<bool> lock;

void lock(){
  while(!(lock->atomic_compare_exchange(&false, true))) { }
}

void unlock(){
  lock = false;
}
\end{lstlisting}
\caption{A simple spinlock implementation of \lstinline{std::mutex}}
\label{mutex-simple}
\end{figure}

We begin by examining \lstinline{std::mutex}, the basic lock class. The class contains methods \lstinline{lock}, \lstinline{unlock}, and \lstinline{try_lock}, as well as a constructor and destructor. The intended behavior is simply to provide mutual exclusion; a simple implementation could be a spinlock as shown in Figure~\ref{mutex-simple}, although more realistic implementations use OS primitives for efficiency. While \lstinline{std::mutex} is the library class corresponding to a standard lock, its informal specification has two restrictions that do not appear in the standard separation logic treatment:
\begin{itemize}
\item When a thread calls \lstinline{unlock}, it must own the mutex (i.e., it must have been the thread that most recently called \lstinline{lock}).
\item When a mutex is destroyed, it must not be owned by any threads.
\end{itemize}
A program that violates either of these conditions has \emph{undefined behavior} according to the library specification, so our program logic must ensure that they are respected.

\begin{mathpar}
  \spectriple
    {\triangleright P}
    {\texttt{mutex::mutex}(m)}
    {\exists\gamma.\;\mathsf{M}_1(m,\gamma,P)\ast\mathsf{token}(\gamma,1)}

  \spectriple
    {\mathsf{M}_1(m,\gamma,P)\ast\mathsf{token}(\gamma,1)}
    {\texttt{mutex::$\sim$mutex}(m)}
    {P}
\end{mathpar}
\begin{mathpar}
  \spectriple
    {\mathsf{M}_q(m,\gamma,P)\ast\mathsf{token}(\gamma,q')}
    {\texttt{lock}(m)}
    {\mathsf{M}_q(m,\gamma,P)\ast P\ast\mathsf{locked}(\gamma,t,q')}

  \spectriple
    {\mathsf{M}_q(m,\gamma,P)\ast\mathsf{locked}(\gamma,t,q')
      \ast\triangleright P}
    {\texttt{unlock}(m)}
    {\mathsf{M}_q(m,\gamma,P)\ast\mathsf{token}(\gamma,q')}

  \spectriple
    {\mathsf{M}_q(m,\gamma,P)\ast\mathsf{token}(\gamma,q')}
    {\texttt{try\_lock}(m)}
    {\begin{gathered}
       \exists b.\;\mathsf{ret}=\mathsf{Vbool}(b)\land
         \mathsf{M}_q(m,\gamma,P)\ast {}\\
       \left(\begin{aligned}
         &b=\mathsf{true}\Rightarrow P\ast\mathsf{locked}(\gamma,t,q')\\
         &b=\mathsf{false}\Rightarrow\mathsf{token}(\gamma,q')
       \end{aligned}\right)
     \end{gathered}}
\end{mathpar}

\subsection{\texorpdfstring{\lstinline{std::recursive_mutex}}{std::recursive\_mutex}}

\begin{mathpar}
  \spectriple
    {\begin{gathered}
       \mathsf{RM}_q(m,\gamma)\ast\mathsf{rmutex}(\gamma,P) \\
       {}\ast\mathsf{acquireable}(\gamma,t,n,P,a)\ast\mathsf{token}(\gamma,q')
     \end{gathered}}
    {\texttt{lock}(m)}
    {\begin{gathered}
       \mathsf{RM}_q(m,\gamma)\ast\mathsf{rmutex}(\gamma,P) \\
       {}\ast\mathsf{acquireable}(\gamma,t,n+1,P,a)\ast\mathsf{given\_token}(\gamma,q')
     \end{gathered}}

  \spectriple
    {\begin{gathered}
       \mathsf{RM}_q(m,\gamma)\ast\mathsf{rmutex}(\gamma,P) \\
       {}\ast\mathsf{acquireable}(\gamma,t,n+1,P,a)\ast\mathsf{given\_token}(\gamma,q')
     \end{gathered}}
    {\texttt{unlock}(m)}
    {\begin{gathered}
       \mathsf{RM}_q(m,\gamma)\ast\mathsf{rmutex}(\gamma,P) \\
       {}\ast\mathsf{acquireable}(\gamma,t,n,P,a)\ast\mathsf{token}(\gamma,q')
     \end{gathered}}
\end{mathpar}

\subsection{\texorpdfstring{\lstinline{std::shared_mutex}}{std::shared\_mutex}}

\begin{mathpar}
  \spectriple
    {\triangleright P(1)}
    {\texttt{shared\_mutex::shared\_mutex}(m)}
    {\exists\gamma.\;\mathsf{SM}_1(m,\gamma,P)\ast
      \mathsf{allUsers}^{\gamma}(\varnothing)}

  \spectriple
    {\mathsf{SM}_1(m,\gamma,P)\ast\mathsf{allUsers}^{\gamma}(\varnothing)}
    {\texttt{shared\_mutex::$\sim$shared\_mutex}(m)}
    {P(1)}
\end{mathpar}
\begin{mathpar}
  \spectriple
    {\mathsf{SM}_q(m,\gamma,P)\ast\mathsf{user}^{\gamma}(t)}
    {\texttt{lock}(m)}
    {\mathsf{SM}_q(m,\gamma,P)\ast P(1)\ast\mathsf{locked}(\gamma,t)}

  \spectriple
    {\mathsf{SM}_q(m,\gamma,P)\ast\mathsf{locked}(\gamma,t)
      \ast\triangleright P(1)}
    {\texttt{unlock}(m)}
    {\mathsf{SM}_q(m,\gamma,P)\ast\mathsf{user}^{\gamma}(t)}

  \spectriple
    {\mathsf{SM}_q(m,\gamma,P)\ast\mathsf{user}^{\gamma}(t)}
    {\texttt{try\_lock}(m)}
    {\begin{gathered}
       \exists b.\;\mathsf{ret}=\mathsf{Vbool}(b)\land
         \mathsf{SM}_q(m,\gamma,P)\ast {}\\
       \left(\begin{aligned}
         &b=\mathsf{true}\Rightarrow P(1)\ast\mathsf{locked}(\gamma,t)\\
         &b=\mathsf{false}\Rightarrow\mathsf{user}^{\gamma}(t)
       \end{aligned}\right)
     \end{gathered}}
\end{mathpar}
\begin{mathpar}
  \spectriple
    {\mathsf{SM}_q(m,\gamma,P)\ast\mathsf{user}^{\gamma}(t)}
    {\texttt{lock\_shared}(m)}
    {\begin{gathered}
       \mathsf{SM}_q(m,\gamma,P)\ast {}\\
       \exists q_P.\;P(q_P)\ast\mathsf{reader\_locked}(\gamma,t,q_P)
     \end{gathered}}

  \spectriple
    {\begin{gathered}
       \mathsf{SM}_q(m,\gamma,P)\ast\mathsf{reader\_locked}(\gamma,t,q_P)\\
       {}\ast\triangleright P(q_P)
     \end{gathered}}
    {\texttt{unlock\_shared}(m)}
    {\mathsf{SM}_q(m,\gamma,P)\ast\mathsf{user}^{\gamma}(t)}

  \spectriple
    {\mathsf{SM}_q(m,\gamma,P)\ast\mathsf{user}^{\gamma}(t)}
    {\texttt{try\_lock\_shared}(m)}
    {\begin{gathered}
       \exists b.\;\mathsf{ret}=\mathsf{Vbool}(b)\land
         \mathsf{SM}_q(m,\gamma,P)\ast {}\\
       \left(\begin{aligned}
         &b=\mathsf{true}\Rightarrow
           \exists q_P.\;P(q_P)\ast\mathsf{reader\_locked}(\gamma,t,q_P)\\
         &b=\mathsf{false}\Rightarrow\mathsf{user}^{\gamma}(t)
       \end{aligned}\right)
     \end{gathered}}
\end{mathpar}

\subsection{The User Library}

TODO cite the Iris reader/writer lock library
TODO mutex is not implemented with the user library yet

\begin{sloppypar}
The \texttt{user} ghost state provides an interface that allows a thread to acquire
locked resources at most once regardless of the mutex kinds. Each mutex
invariant is associated with a disjoint set of thread IDs, managed by one
authoritative piece \texttt{allUsers} and fragment pieces \texttt{user}:
\begin{align*}
  \mathsf{allUsers}^{\gamma}(S)
    &\eqdef \mathsf{own}^{\gamma}(\mathord{\bullet}S), \\
  % \mathsf{users}^{\gamma}(T)
  %   &\eqdef \mathsf{own}^{\gamma}(\mathord{\circ}T), \\
  \mathsf{user}^{\gamma}(t)
    &\eqdef \mathsf{own}^{\gamma}(\mathord{\circ}\{t\}).
\end{align*}
The \texttt{allUsers($S$)} ghost state tracks all "logged-in" threads $S$. Each thread with thread ID $t$ can get
a \texttt{user(t)} handle for that mutex with the \texttt{login} rule, and the handle is
unique for each thread (user-unique). The mutex invariants are slightly
different for different kinds of mutexes, but they all ensure that a thread
must trade in its unique \texttt{user} handle to acquire mutex protected resources.
This ensures that a thread can't acquire protected resources if it already has
some, and unsafe programs that calls \texttt{mutex.lock()} or
\texttt{shared\_mutex.lock\_shared()} twice consecutively in a single thread
would get stuck. Notably, for recursive mutexes, if a thread already has the
resource, it is safe to \texttt{lock()} again, and the thread simply does not
exchange anything with the invariant.

A thread can also "logout" and turn in its \texttt{user} handle. Deallocation
of any kind of mutex requires no "logged-in" users
($\mathsf{allUsers}^{\gamma}(\varnothing)$), enforcing that any allocated
\texttt{user} handle is deallocated with \texttt{logout}
(\texttt{all-user-logged-out}), and so the mutex is not held
by any thread, thus safe to deallocate.


\begin{figure}[t]
\begin{equation}
  \mathsf{user}^{\gamma}(t)\ast\mathsf{user}^{\gamma}(t)
  \vdash \mathsf{False}
  \tag{\textsc{user-unique}}
  \label{law:user-unique}
\end{equation}
\begin{equation}
  t\notin S \rightarrow
  \mathsf{allUsers}^{\gamma}(S)
  \vdash \upd
  \mathsf{allUsers}^{\gamma}(\{t\}\cup S)
  \ast\mathsf{user}^{\gamma}(t)
  \tag{\textsc{login}}
  \label{law:user-login}
\end{equation}
\begin{equation}
  t\notin S \rightarrow
  \mathsf{allUsers}^{\gamma}(\{t\}\cup S)\ast\mathsf{user}^{\gamma}(t)
  \vdash \upd \mathsf{allUsers}^{\gamma}(S)
  \tag{\textsc{logout}}
  \label{law:user-logout}
\end{equation}
\begin{equation}
  \mathsf{allUsers}^{\gamma}(\varnothing)\ast\mathsf{user}^{\gamma}(t)
  \vdash \mathsf{False}
  \tag{\textsc{all-user-logged-out}}
  \label{law:all-user-logged-out}
\end{equation}
\caption{Rules for the user ghost state.}
\label{fig:user-rules}
\end{figure}


\end{sloppypar}
